\subsection{Resistive Voltage Divider}

The block is constructed from two resistors in series and is used
to attenuate a voltage signal. It is one of the simplest and most
commonly used circuit in electronics.

\subsubsection{Schematic}
\begin{circuitfig}
	% Paths, nodes and wires:
	\node[ground, rotate=-90] at (5.5, 5.25){};
	\node[ground, rotate=-90] at (5.5, 4.75){};
	\draw (5.5, 6.75) to[sinusoidal voltage source, l={$V_{cc}$}] (5.5, 5.25);
	\draw (5.5, 4.75) to[sinusoidal voltage source, l={$V_{ee}$}] (5.5, 3.25);
	\draw (7, 5.25) to[european resistor, l_={$G_{1}$}] (7, 7);
	\draw (7, 3) to[european resistor, l_={$G_{2}$}] (7, 4.75);
	\draw (7, 5) -- (7.75, 5);
	\draw (7, 7) -| (5.5, 7) -| (5.5, 6.75);
	\node[shape=rectangle, minimum width=0.715cm, minimum height=0.465cm](N1) at (8.125, 5){} node[anchor=center] at (N1.text){$V_{out}$};
	\draw (7, 5.25) -| (7, 5);
	\draw (7, 4.75) -| (7, 5);
	\draw (5.5, 3.25) -| (5.5, 3) -- (7, 3);
\end{circuitfig}

\subsubsection{Transfer Function}
Applying KCL at $V_{out}$ we get:
\[
  (V_{out} - V_{cc}){G_1} + (V_{out} - V_{ee}){G_2} = 0
\]
\paragraph{Superposition:}
\begin{enumerate}
  \item Due to \(V_{cc}\) alone (\(V_{ee} = 0V\)):
    \begin{align*}
      (V_{out} - V_{cc}){G_1} + (V_{out} - 0){G_2} &= 0 \\
      V_{out}G_1 - V_{cc}G_1 + V_{out}G_2 &= 0 \\
      V_{out}(G_1 + G_2) - V_{cc}G_1 &= 0 \\
      V_{out}(G_1 + G_2) &= V_{cc}G_1 \\
      H_1(s) = \frac{N_1(s)}{D_1(s)} = \frac{V_{out_1}}{V_{cc}} &= \frac{G_1}{G_1 + G_2}
    \end{align*}
  \item Due to \(V_{ee}\) alone (\(V_{cc} = 0V\)):
    \begin{align*}
      (V_{out} - 0){G_1} + (V_{out} - V_{ee}){G_2} &= 0 \\
      V_{out}G_1 + V_{out}G_2 - V_{ee}G_2 &= 0 \\
      V_{out}(G_1 + G_2) - V_{ee}G_2 &= 0 \\
      V_{out}(G_1 + G_2) &= V_{ee}G_2 \\
      H_2(s) = \frac{N_2(s)}{D_2(s)} = \frac{V_{out_2}}{V_{ee}} &= \frac{G_2}{G_1 + G_2}
    \end{align*}
\end{enumerate}
\paragraph{Total Output:}
\begin{align*}
  V_{out} &= V_{out_1} + V_{out_2} \\
  V_{out} &= V_{cc} \frac{G_1}{G_1 + G_2} + V_{ee} \frac{G_2}{G_1 + G_2}
\end{align*}

\subsubsection{Analysis}
\begin{enumerate}
  \item \textbf{Resistor Ratio (\(\frac{G_1}{G_2}\)):}
    From KCL at \(V_{out}\):
    \begin{align*}
      (V_{out} - V_{cc})G_1 - (V_{ee} - V_{out})G_2 &= 0 \\
      (V_{out} - V_{cc})G_1 &= (V_{ee} - V_{out})G_2 \\
      \frac{G_1}{G_2} &= \frac{V_{ee} - V_{out}}{V_{out} - V_{cc}} \\
      \frac{R_2}{R_1} &= \frac{V_{ee} - V_{out}}{V_{out} - V_{cc}} \hspace{1cm} \text{(Substituting \(G_1 = \frac{1}{R_1}\) and \(G_2 = \frac{1}{R_2}\))}
    \end{align*}
  \item \textbf{Impedance's:}
    \paragraph{At Input:}
      \begin{align*}
        Z_{in} &= \frac{1}{G_1} + \frac{1}{G_2} \\
        Z_{in} &= R_1 + R_2 \hspace{1cm} \text{(Substituting \(G_1 = \frac{1}{R_1}\) and \(G_2 = \frac{1}{R_2}\))}
      \end{align*}
    \paragraph{At Output:}
      \begin{align*}
        Z_{out} &= \frac{1}{G_1 + G2} \\
        Z_{out} &= \frac{1}{\frac{1}{R_1} + \frac{1}{R_2}} \hspace{1cm} \text{(Substituting \(G_1 = \frac{1}{R_1}\) and \(G_2 = \frac{1}{R_2}\))} \\
        Z_{out} &= \frac{1}{\frac{R_1 + R_2}{R_1 R_2}} \\
        Z_{out} &= \frac{R_1 R_2}{R_1 + R_2}
      \end{align*}
  \item \textbf{Loading Effect:}
    Consider a load impedance \(Z_L\) connected to the output of the voltage
    divider. This impedance will itself form a voltage divider with our
    resistive divider. It is thus recommended to keep the output impedance
    of the divider at least 10 times smaller than the load impedance at the
    respective frequency of operation or If you can control the load,
    keep the load impedance at least 10 times larger than the output
    impedance of the divider.
\end{enumerate}

\subsubsection{Question}
Design a resistive voltage divider to step down a \(10V\) supply to
\(3.3V\). Calculate the input and output impedance of the voltage divider.
\paragraph{Solution:}
\begin{enumerate}
  \item Calculate the resistor ratio (\(\frac{R_2}{R_1}\)):
    \begin{align*}
      \frac{R_2}{R_1} &= \frac{V_{ee} - V_{out}}{V_{out} - V_{cc}} \\
      &= \frac{0V - 3.3V}{3.3V - 10V} \\
      &= \frac{-3.3V}{-6.7V} \\
    \frac{R_2}{R_1} &\approx 0.4925
    \end{align*}
  \item Fix one of the resistors (\(R_2 = 10 k\Omega\)) to a standard
    value and calculate the other resistor (\(R_1\)):
    \begin{align*}
      R_1 &= \frac{R_2}{0.4925} \\
      &= \frac{10 k\Omega}{0.4925} \\
      R_1 &\approx 20.3 k\Omega
    \end{align*}
  \item Calculating Input and Output Impedances:
    \[
        \begin{aligned}
        Z_{in} &= R_1 + R_2 \\
        &= 20.3 k\Omega + 10 k\Omega \\
        Z_{in} &\approx 30.3 k\Omega
        \end{aligned}
      \]
      \[
        \begin{aligned}
          Z_{out} &= \frac{R_1 R_2}{R_1 + R_2} \\
          Z_{out} &= \frac{(20.3 k\Omega)(10 k\Omega)}{20.3 k\Omega + 10 k\Omega} \\
          Z_{out} &\approx 6.77 k\Omega
        \end{aligned}
      \]
\end{enumerate}

